\documentclass[journal]{IEEEtran}

\usepackage{amsmath,amssymb,amsthm}
\usepackage{kotex}
\usepackage{booktabs}
\usepackage{array}
\usepackage{tabularx}
\usepackage{tikz}
\usetikzlibrary{arrows.meta,positioning,shapes.geometric,fit,backgrounds}
\usepackage{xurl}
\usepackage[hidelinks,unicode]{hyperref}
\usepackage{cite}

% Portable bundled TeX Live font setup (Times-equivalent + Hangul Un-fonts + Unicode Mono)
\setmainfont{TeXGyreTermesX}
\setsansfont{texgyreheros-regular.otf}[
  BoldFont = texgyreheros-bold.otf,
  ItalicFont = texgyreheros-italic.otf,
  BoldItalicFont = texgyreheros-bolditalic.otf,
  Scale = MatchLowercase
]
\setmonofont{DejaVuSansMono.ttf}[
  BoldFont = DejaVuSansMono-Bold.ttf,
  ItalicFont = DejaVuSansMono-Oblique.ttf,
  BoldItalicFont = DejaVuSansMono-BoldOblique.ttf,
  Scale = 0.80
]
\setmainhangulfont{UnBatang.ttf}[
  BoldFont = UnBatangBold.ttf,
  BoldItalicFont = UnBatangBold.ttf,
  BoldItalicFeatures = {FakeSlant=0.15},
  AutoFakeSlant = 0.15
]
\setsanshangulfont{UnDotum.ttf}[
  BoldFont = UnDotumBold.ttf,
  BoldItalicFont = UnDotumBold.ttf,
  BoldItalicFeatures = {FakeSlant=0.15},
  AutoFakeSlant = 0.15
]

% Preserve IEEEtran visual hierarchy while preventing Hangul spacing collapse under \scshape
\makeatletter
\long\def\@makecaption#1#2{%
  \ifx\@captype\@IEEEtablestring%
    \footnotesize\bgroup\par\centering\@IEEEtabletopskipstrut{\normalfont\footnotesize #1}\\{\normalfont\footnotesize\bfseries #2}\par\addvspace{0.5\baselineskip}\egroup%
    \@IEEEtablecaptionsepspace
  \else%
    \@IEEEfigurecaptionsepspace
    \setbox\@tempboxa\hbox{\normalfont\footnotesize {#1.}\nobreakspace\nobreakspace #2}%
    \ifdim \wd\@tempboxa >\hsize%
      \setbox\@tempboxa\hbox{\normalfont\footnotesize {#1.}\nobreakspace\nobreakspace}%
      \parbox[t]{\hsize}{\normalfont\footnotesize\noindent\unhbox\@tempboxa#2}%
    \else%
      \ifCLASSOPTIONconference \hbox to\hsize{\normalfont\footnotesize\hfil\box\@tempboxa\hfil}%
      \else \hbox to\hsize{\normalfont\footnotesize\box\@tempboxa\hfil}%
      \fi%
    \fi%
  \fi%
}
\def\section{\@startsection{section}{1}{\z@}{1.5ex plus 1.5ex minus 0.5ex}%
  {0.7ex plus 1ex minus 0ex}{\normalfont\normalsize\centering\bfseries}}%
\makeatother

\emergencystretch 1.8em
\tracinglostchars=2
\setlength{\textfloatsep}{10pt plus 2pt minus 2pt}
\setlength{\floatsep}{8pt plus 2pt minus 2pt}
\setlength{\intextsep}{8pt plus 2pt minus 2pt}

\newcommand{\code}[1]{\texttt{#1}}
\newtheorem{definition}{Definition}

\begin{document}

\title{Proof-Notebook: Lean 4 및 Rocq 다중 증명기 노트북 커널과 협업 웹 스튜디오 아키텍처}

\author{이인수%
\thanks{본 기술 보고서는 IRADE Proof-Notebook 프로젝트(\code{notebook-kernel}, 패키지명 \code{lean4-jupyter-kernel} v0.1.0, Apache-2.0)의 아키텍처, 프로토콜 구현 및 검증 결과를 기록한다.}}

\markboth{IRADE Proof-Notebook Technical Report, 2026}{이인수: Lean 4 및 Rocq 다중 증명기 노트북 아키텍처}

\maketitle

\begin{abstract}
대화형 정리 증명기(Interactive Theorem Prover, ITP)는 형식 검증과 수학적 추론에 필수적이나, 선형 셀 실행을 가정하는 전통적 Jupyter 노트북 환경에 통합할 경우 상태 비일관성, 심볼 재선언 충돌, 그리고 커서 위치 기반 목표(Goal) 상태 가시성 부재라는 구조적 문제가 발생한다. 본 보고서는 Lean 4(\code{lean4\_kernel}), Rocq/Coq(\code{rocq\_kernel}), 및 Python 3(\code{ipykernel})을 단일 워크벤치 아래 결합한 \textit{Proof-Notebook} 시스템(저장소 내 총 17,183라인 규모)의 설계와 구현을 기술한다. 제안 시스템은 (1) \code{leanprover-community/repl}의 JSON 스냅샷 ID 기반 불변 환경 방향성 비순환 그래프(DAG)와 자동 백트래킹 엔진, (2) \code{coqidetop}의 네이티브 XML 상태 트랜잭션 머신(STM) 기반 어휘 토크나이징 및 셀 단위 원자적 롤백(\code{Edit\_at}) 엔진, (3) \code{InfoTree} 및 STM 스텝을 커서와 동기화하는 하이브리드 서버 사이드/클라이언트 사이드 렌더링(SSR+CSR) InfoView 위젯, 그리고 (4) TLS 1.2+ HTTPS 전용 서버, 제로 의존성 RESP2 Redis 핫/콜드 워크스페이스 생명주기 관리, SSE 실시간 협업 및 Project-Agora 캔버스 브리지를 갖춘 VS Code 스타일 웹 스튜디오로 구성된다. 전체 30개 자동화 테스트 중 오프라인 프로토콜 및 라이브 Rocq 세션을 포함한 21개 항목이 통과하였으며, 샌드박스 환경의 \code{lean} 바이너리 부재(2건 스킵)와 루프백 프록시 제한(7건 실패) 및 웹 프론트엔드 내 Rocq UI 노출 범위를 포함한 현 통합 경계까지 투명하게 보고한다.
\end{abstract}

\begin{IEEEkeywords}
Lean 4, Rocq, Coq, Jupyter 커널, 대화형 정리 증명, InfoView, 스냅샷 DAG, 상태 트랜잭션 머신(STM), 실시간 협업 노트북.
\end{IEEEkeywords}

\section{서론}

\subsection{배경 및 연구 동기}
계산 과학 워크플로우에서 Jupyter 노트북은 서술형 텍스트(Markdown), 수식(\LaTeX), 실행 코드 및 풍부한 시각 출력을 단일 문서로 결합하여 재현 가능한 연구 매체로 자리 잡았다~\cite{kluyver2016jupyter}. 한편, Lean 4~\cite{moura2021lean4}와 Rocq(구 Coq)~\cite{barras2015asynchronous}와 같은 의존 타입 이론(Dependent Type Theory) 기반 대화형 정리 증명기는 프로그램 검증과 형식 수학에서 빠르게 확산되고 있으나, 이를 노트북 패러다임으로 옮기는 과정은 단순한 스크립트 인터프리터 래핑과 근본적으로 다르다.

기존의 문학적 증명(Literate Proof) 도구인 Alectryon~\cite{pitclaudel2020alectryon}은 증명 스크립트와 중간 목표 상태를 정적 HTML로 엮어 가독성을 높였으나, 연구자가 셀을 임의 순서로 수정·재실행하고 파이썬 수치/기호 실험과 형식 증명을 동일 작업 공간에서 실시간 협업으로 연결하는 대화형 환경으로는 확장되지 않았다. 본 프로젝트(\code{proof-notebook}, \code{pyproject.toml} 기준 \code{lean4-jupyter-kernel} v0.1.0)는 표준 Jupyter 커널 사양과 VS Code 스타일의 독립 협업 웹 스튜디오(\textit{IRADE Proof Studio})를 동시에 제공함으로써 이러한 공백을 해소하도록 설계되었다.

\subsection{핵심 기술적 난제}
다중 증명기 노트북 시스템을 구축할 때 해결해야 할 핵심 난제는 크게 네 가지이다.
\begin{enumerate}
  \item \textbf{비선형 셀 재실행과 전역 이름 충돌:} Jupyter 사용자는 상단 셀의 \code{def}나 \code{theorem}을 수정한 뒤 하단 셀을 다시 실행한다. 단순 누적 프로세스는 동일 식별자에 대해 \code{already been declared} 오류를 유발하므로, 과거 환경 스냅샷으로 되돌려 후속 계보(Lineage)를 재구성해야 한다.
  \item \textbf{하위 프로세스 중단 및 복구:} 타임아웃, 비정상 종료 또는 사용자 인터럽트 발생 시 장기 실행 중인 Lean/Rocq 프로세스 상태를 잃지 않고 유효한 증명 계보를 자동 복원해야 한다.
  \item \textbf{행/열 커서 기반 InfoView 제공:} 증명기는 최종 성공/실패 메시지보다 각 전술(Tactic) 경계 전후의 가설 맥락($\Gamma$)과 목표($\vdash G$) 변화가 핵심이다. 커널 실행 결과를 정적 문자열이 아닌 커서 연동 위젯으로 전달해야 한다.
  \item \textbf{이기종 증명기 프로토콜의 통합:} Lean 4 REPL은 빈 줄(\code{\textbackslash n\textbackslash n})로 구분된 비동기 JSON 스냅샷 ID 모델~\cite{morrison_repl}을 사용하고, Rocq \code{coqidetop}은 XML 직렬화 기반의 STM(\code{Init}, \code{Add}, \code{Edit\_at}, \code{Goal})~\cite{barras2015asynchronous}을 사용한다. 상이한 두 상태 모델을 공통된 노트북 사용자 경험으로 추상화하면서도 각 증명기의 안전성 보장을 훼손하지 않아야 한다.
\end{enumerate}

\subsection{주요 구현 기여 및 보고서 구성}
본 보고서가 기술하는 시스템의 주요 엔지니어링 기여는 다음과 같다.
\begin{itemize}
  \item \textbf{Lean 4 불변 스냅샷 DAG 커널(\code{lean4\_kernel}):} \code{EnvNode}/\code{ProofNode} 계보 추적, 재선언 자동 백트래킹, \code{InfoTree} 재귀 파싱 기반 하이브리드 InfoView, 17개 매직 명령 및 172개 유니코드 약어 완성.
  \item \textbf{Rocq XML STM 커널(\code{rocq\_kernel}):} 외부 OPAM 패키지 추가 없이 표준 \code{coqidetop}과 직접 통신하는 XML 드라이버, 중첩 주석과 줄임표(\code{..})를 인지하는 어휘 문장 토크나이저, 그리고 오류 발생 시 셀 전체를 진입 상태로 되돌리는 원자적 롤백.
  \item \textbf{다중 커널 Web Proof Studio 및 협업 인프라(\code{web.py}, \code{static/}):} LRU 16슬롯 ZMQ 커널 풀, 외부 라이브러리 없는 RESP2+TLS Redis 클라이언트와 HOT/COLD 생명주기 관리자, SSE 실시간 동기화 및 Project-Agora 양방향 캔버스 브리지.
\end{itemize}

이후 제2장에서는 전체 4계층 아키텍처와 코드베이스 규모를 개관하고, 제3장과 제4장에서는 각각 Lean 4 커널과 Rocq 커널의 프로토콜·상태 관리·InfoView 메커니즘을 상세히 분석한다. 제5장은 웹 플랫폼과 협업 생명주기를 다루며, 제6장에서는 CLI 진단·예제 증명·테스트 스위트 검증 결과와 함께 현 통합상의 공학적 한계를 평가하고, 제7장에서 결론을 맺는다.

\section{전체 시스템 아키텍처 개요}

\subsection{이중 실행 경로 기반 4계층 아키텍처}

\begin{figure*}[!t]
\centering
\begin{tikzpicture}[
  font=\footnotesize,
  box/.style={draw, rounded corners=2pt, align=center, inner sep=4pt, minimum height=9mm},
  layer/.style={draw=black!45, dashed, rounded corners=3pt, inner sep=9pt},
  layerlbl/.style={fill=white, inner sep=1.5pt, font=\scriptsize\bfseries, anchor=west},
  arr/.style={-{Stealth[length=2.2mm]}, semithick},
  biarr/.style={{Stealth[length=2.0mm]}-{Stealth[length=2.0mm]}, semithick}
]
% Layer 1: Clients
\node[box, fill=blue!6, text width=4.3cm] (jlab) {JupyterLab / nbclient / IDE\\(표준 Jupyter ZMQ 클라이언트)};
\node[box, fill=blue!6, text width=9.9cm, right=0.8cm of jlab] (webclient) {VS Code 스타일 Web Proof Studio (\code{index.html}, \code{app.js}, \code{styles.css})\\Explorer · Editor Tabs · SSE Sync · InfoView (\code{mountInteractiveHtml}) · Agora Bridge};

% Layer 2: Transport & Control Plane
\node[box, fill=teal!6, text width=4.3cm, below=0.75cm of jlab] (zmq) {표준 Jupyter ZMQ 프로토콜\\(\code{ipykernel.kernelbase.Kernel})};
\node[box, fill=teal!6, text width=9.9cm, below=0.75cm of webclient] (https) {TLS 1.2+ HTTPS 서버 (\code{ProofNotebookHTTPServer} / \code{ProofStudioApp})\\\code{ProofKernelPool} (LRU $\le 16$) · \code{WorkspaceLifecycleManager} · \code{RealtimeNotebookHub} (SSE)};

% Layer 3: Kernel Engines
\node[box, fill=yellow!10, text width=4.5cm, below=0.95cm of https.south west, anchor=north west, xshift=-3.6cm] (klean) {\code{lean4\_kernel} (\code{Lean4Kernel})\\\code{LeanSession} (DAG) · \code{infoview} · 17 Magics};
\node[box, fill=yellow!10, text width=4.5cm, right=0.45cm of klean] (krocq) {\code{rocq\_kernel} (\code{RocqKernel})\\\code{RocqSession} (STM) · Tokenizer · 8 Magics};
\node[box, fill=yellow!10, text width=3.5cm, right=0.45cm of krocq] (kpy) {\code{ipykernel\_launcher}\\Python 3 연산 커널};

% Layer 4: Prover Subprocesses & Persistence
\node[box, fill=gray!10, text width=4.5cm, below=0.80cm of klean] (plean) {\code{REPLProcess} (\code{repl} / \code{lake env repl})\\JSON + \code{\textbackslash n\textbackslash n} 프레이밍 (\code{v4.35.0})};
\node[box, fill=gray!10, text width=4.5cm, below=0.80cm of krocq] (procq) {\code{RocqProcess} (\code{coqidetop.opt})\\XML + \code{-main-channel stdfds}};
\node[box, fill=green!7, text width=3.5cm, below=0.80cm of kpy] (store) {\code{AgoraRedisClient} (HOT, 3600s)\\\code{NotebookStore} (Disk COLD)\\\code{AgoraBridge} (HTTPX / Canvas)};

% Arrows
\draw[biarr] (jlab) -- node[right, font=\scriptsize]{ZMQ} (zmq);
\draw[biarr] (webclient) -- node[right, font=\scriptsize]{HTTPS / SSE} (https);
\draw[biarr] (zmq.south) -- (klean.north -| zmq.south);
\draw[arr] (https.south -| klean.north) -- node[pos=0.68, left, font=\scriptsize]{\code{lean4}} (klean.north);
\draw[arr] (https.south -| krocq.north) -- node[pos=0.68, left, font=\scriptsize]{\code{rocq}} (krocq.north);
\draw[arr] (https.south -| kpy.north) -- node[pos=0.68, right, font=\scriptsize]{\code{python}} (kpy.north);
\draw[biarr] (klean) -- node[pos=0.60, right, font=\scriptsize]{stdin/out} (plean);
\draw[biarr] (krocq) -- node[pos=0.60, right, font=\scriptsize]{stdfds} (procq);
\draw[biarr] (https.east) -- ++(0.65,0) |- node[pos=0.25, right, font=\scriptsize, align=left]{HOT/COLD\\동기화} (store.east);

% Group background labels
\begin{scope}[on background layer]
  \node[layer, fit=(jlab)(webclient)] (l1) {};
  \node[layer, fit=(zmq)(https)] (l2) {};
  \node[layer, fit=(klean)(krocq)(kpy)] (l3) {};
  \node[layer, fit=(plean)(procq)(store)] (l4) {};
  \node[layerlbl] at ([xshift=4mm]l1.north west) {1. Presentation};
  \node[layerlbl] at ([xshift=4mm]l2.north west) {2. Transport / Pool};
  \node[layerlbl, anchor=center] at ([xshift=-2mm]kpy.west |- l3.north) {3. Kernel Logic};
  \node[layerlbl, anchor=center] at ([xshift=-2mm]store.west |- l4.north) {4. Prover \& State};
\end{scope}
\end{tikzpicture}
\caption{Proof-Notebook 4계층 하이브리드 아키텍처. 표준 Jupyter ZMQ 실행 경로와 HTTPS 기반 Web Proof Studio 경로가 동일한 커널 엔진(\code{lean4\_kernel}, \code{rocq\_kernel}, \code{ipykernel})을 공유한다.}
\label{fig:architecture}
\end{figure*}

Proof-Notebook 시스템은 Fig.~\ref{fig:architecture}에 나타낸 바와 같이 네 개의 계층과 두 개의 접속 경로로 구조화된다. 첫 번째 경로는 표준 Jupyter ZeroMQ 프로토콜을 통해 JupyterLab, VS Code Notebook 또는 \code{nbclient}가 커널 스펙(\code{lean4}, \code{rocq})에 직접 연결되는 전통적 방식이다. 두 번째 경로는 \code{lean4-proof-studio}(또는 \code{lean4-kernel serve}) 명령으로 기동하는 전용 웹 워크벤치로서, 단일 워크스페이스 안에 Lean 4(및 백엔드 풀의 Rocq), Python 3, Markdown 셀을 혼합 배치하고 Redis·디스크·Project-Agora 플랫폼 간 양방향 동기화를 수행한다.

\code{pyproject.toml}은 Python 3.10 이상을 요구하며 핵심 런타임 의존성을 \code{ipykernel>=6.20}, \code{jupyter\_client>=8.0}, \code{httpx>=0.25} 단 3종으로 최소화하였다(\code{requirements.txt} 기준 \code{ipykernel==7.4.0}, \code{jupyter\_client==8.10.0}, \code{pyzmq==27.3.0} 등에서 검증). 외부 Redis 클라이언트 라이브러리조차 내장 RESP2 구현(\code{AgoraRedisClient})으로 대체하여 배포 복잡도를 낮췄다.

\subsection{모듈 구성 및 코드베이스 규모}
저장소의 실측 코드 라인 수(\code{wc -l})와 주요 모듈 책임은 Table~\ref{tab:codebase_metrics}에 정리하였다. Python 커널과 웹 서버(\code{src/lean4\_kernel/} 8,032라인, \code{src/rocq\_kernel/} 2,949라인), 정적 프론트엔드(\code{static/} 4,403코드라인), 테스트 스위트(1,578라인) 및 예제(66코드라인, 문서 포함 151라인)를 합산하면 총 17,183라인에 달한다. 특히 \code{src/lean4\_kernel/web.py}(3,395라인)와 \code{app.js}(2,338라인)가 협업 스튜디오의 제어 평면과 UI 평면을 담당한다.

\begin{table}[!t]
\caption{Proof-Notebook 모듈별 실측 규모 (\code{wc -l}) 및 핵심 역할}
\label{tab:codebase_metrics}
\centering
\scriptsize
\begin{tabularx}{\columnwidth}{@{}l r X@{}}
\toprule
\textbf{구성 요소 / 경로} & \textbf{라인 수} & \textbf{주요 클래스 및 기능} \\
\midrule
\multicolumn{3}{@{}l}{\textbf{Lean 4 커널 패키지 (\code{src/lean4\_kernel/}, 소계 8,032라인)}} \\
\code{repl.py} & 601 & \code{REPLBinaryManager}, \code{REPLProcess}, JSON 프레이밍 \\
\code{session.py} & 974 & \code{LeanSession}, \code{EnvNode}/\code{ProofNode} DAG, 복구 \\
\code{infoview.py} & 1,260 & \code{InfoTree} 추출, SSR+CSR split-pane 위젯 \\
\code{magics.py} & 454 & 17종 라인/셀 매직 디스패처 (\code{dispatch\_magic}) \\
\code{renderer.py} & 309 & \code{text/plain}·\code{text/html} 이중 MIME 렌더러 \\
\code{kernel.py} & 440 & \code{Lean4Kernel} (\code{do\_execute}/\code{complete}/\code{inspect}) \\
\code{unicode\_map.py} & 311 & 유니코드 172, 키워드 62, 전술 56 매핑 \\
\code{install.py} & 264 & 커널스펙 설치·\code{doctor}·CLI 서브커맨드 \\
\code{web.py} & 3,395 & HTTPS 스튜디오, 커널 풀, Redis/Agora 브리지 \\
\code{\_\_init\_\_.py}, \code{\_\_main\_\_.py} & 24 & 패키지 메타데이터(\code{0.1.0}) 및 진입점 \\
\midrule
\multicolumn{3}{@{}l}{\textbf{Rocq 커널 패키지 (\code{src/rocq\_kernel/}, 소계 2,949라인)}} \\
\code{xml\_protocol.py} & 395 & XML 빌더/파서, 스트림 위생 처리 \\
\code{driver.py} & 288 & \code{RocqProcess} (\code{coqidetop -main-channel stdfds}) \\
\code{session.py} & 689 & \code{tokenize\_sentences}, \code{STMNode}, 원자적 롤백 \\
\code{infoview.py} & 685 & \code{RocqStepInfo} 추출, split-pane InfoView \\
\code{kernel.py} & 420 & \code{RocqKernel}, 8종 매직, 완결성 검사 \\
\code{unicode\_map.py} & 283 & 유니코드 123, 키워드 66, 전술 74 매핑 \\
\code{install.py} & 159 & Rocq 커널스펙 설치 및 \code{Check nat.} 진단 \\
\code{\_\_init\_\_.py}, \code{\_\_main\_\_.py} & 30 & 패키지 공개 API 및 CLI 진입점 \\
\midrule
\multicolumn{3}{@{}l}{\textbf{웹 프론트엔드, 테스트, 예제 및 데모 노트북}} \\
\code{static/\{html,js,css,md\}} & 4,473 & \code{app.js}(2,338), \code{styles.css}(1,748), \code{index.html}(317), \code{README}(70) \\
\code{tests/*.py} & 1,578 & 6개 테스트 파일, 총 30개 테스트 케이스 \\
\code{examples/*} & 151 & \code{Arithmetic.lean}(23), \code{Arithmetic.v}(43), \code{README}(85) \\
\midrule
\textbf{측정 합계 (코드/테스트)} & \textbf{17,183} & 별도 \code{notebooks/*.ipynb}(2종), \code{vendor/repl/} 포함 \\
\bottomrule
\end{tabularx}
\end{table}

\section{Lean 4 커널 엔진 (\code{lean4\_kernel})}

\subsection{REPL 프로세스 수명주기 및 JSON 프레이밍 (\code{repl.py})}
\code{lean4\_kernel}은 저장소 내 \code{vendor/repl/}(\code{leanprover/lean4:v4.35.0}, \code{lakefile.toml}에서 \code{supportInterpreter = true} 활성화)에 포함된 \code{leanprover-community/repl} 바이너리~\cite{morrison_repl}를 하위 프로세스로 구동한다. \code{REPLBinaryManager}는 우선 프로젝트 로컬 또는 \code{vendor/repl} 빌드를 탐색하고, 부재 시 툴체인 버전별로 \code{\textasciitilde/.cache/lean4-kernel/repl} 경로에 컴파일된 캐시를 관리한다. \code{resolve\_lean\_sysroot} 및 \code{detect\_lean\_version}은 \code{elan} 기반 작업 공간의 프로젝트 컨텍스트를 판별하며, 유효한 \code{project\_dir}(예: \code{lakefile.toml}/\code{lean-toolchain})가 감지되면 \code{REPLProcess}는 \code{lake env repl}을 실행하여 해당 Mathlib/프로젝트 의존성을 상속한다.

프로세스 간 통신(IPC)은 \code{stdin}/\code{stdout} 상에서 공백 줄(\code{\textbackslash n\textbackslash n})로 구분된 단일 라인 콤팩트 JSON 메시지를 교환한다. 모든 요청은 \code{\_req\_lock}으로 직렬화되며, 백그라운드 데몬 스레드(\code{\_stdout\_reader}, \code{\_stderr\_reader})가 파이프 교착 상태 없이 비동기 출력을 큐에 적재한다. 특히 Lean C 런타임이나 \code{IO.println}이 JSON 블록 직전에 원시 표준 출력을 방출하는 경우 \code{parse\_repl\_json\_block}이 선행 텍스트를 \code{aux\_stdout}으로 분리 추출하고 뒤따르는 JSON 객체만 안전하게 역직렬화한다. 종료 및 타임아웃 복구 시에는 \code{start\_new\_session=True}로 생성된 프로세스 그룹 전체에 \code{SIGTERM}을 전송한 뒤 1초 유예 후 \code{SIGKILL}을 발행하여 고아 Lean 프로세스 잔류를 차단한다.

\subsection{불변 스냅샷 DAG와 자동 백트래킹 (\code{session.py})}
\code{REPLProcess}가 반환하는 정수형 \code{env} 및 \code{proofState} 핸들은 부수 효과 없는 불변 스냅샷을 가리킨다. \code{LeanSession}은 이를 명시적 방향성 비순환 그래프(DAG)로 모델링한다.

\begin{definition}[환경 스냅샷 DAG]
Lean 세션의 명령 실행 이력을 $G_E = (V_E, E_E)$라 할 때, 각 노드 $v \in V_E$는 \code{EnvNode} 튜플 $(e_v, p_v, c_v, \text{cid}_v, \mathcal{D}_v)$로 정의된다. 여기서 $e_v \in \mathbb{N}_{\ge 0}$는 REPL 환경 ID, $p_v \in V_E \cup \{\bot\}$는 부모 노드(\code{parent\_env}), $c_v$는 정규화된 소스 코드, $\text{cid}_v$는 Jupyter 셀 식별자, $\mathcal{D}_v$는 \code{extract\_declared\_names}가 추출한 선언 심볼 집합이다. 현재 활성 노드를 $v_{\text{curr}}$라 하면 루트 $v_0$에서 $v_{\text{curr}}$까지의 유일한 계보(Lineage)는 $\pi(v_{\text{curr}}) = (v_0, v_1, \dots, v_k)$이다.
\end{definition}

전술 모드 상태 역시 \code{ProofNode}(\code{proof\_state}, \code{parent\_proof}, \code{env\_id}, \code{tactic})로 분리 추적되어 \code{\%tactic} 또는 \code{\%\%tactic}을 통한 다단계 증명 탐색과 \code{sorry} 기반 상호작용을 지원한다.

노트북 셀을 비순차적으로 재실행할 때 발생하는 대표적 오류는 동일 이름의 재선언이다. \code{LeanSession.\_try\_auto\_backtrack}은 다음 절차로 이를 자동 해결한다.
\begin{enumerate}
  \item REPL 오류 메시지에서 유니코드 곡선 따옴표(\code{‘...’})와 ASCII 따옴표(\code{'...'})를 모두 정규 표현식으로 매칭하여 충돌 식별자 $\sigma$ 또는 동일 \code{cell\_id}를 탐색한다.
  \item \code{active\_lineage()}를 따라 상위 조상을 역추적하여 $\sigma \in \mathcal{D}_{v_j}$ 또는 $\text{cid}_{v_j} = \text{cid}_{\text{req}}$를 만족하는 가장 이른 충돌 노드 $v_j$를 특정한다.
  \item 활성 포인터를 $p_{v_j}$(\code{parent\_env})로 되돌리고, 수정된 새 셀을 실행하여 분기 노드 $v_j'$를 생성한 뒤 기존 후손 계보 $(v_{j+1}, \dots, v_k)$를 순차 재실행한다. 단, 헤더 \code{import}가 선두에 존재하면(\code{starts\_with\_header\_import}) REPL의 최상위 \code{env=None} 초기화 규약에 맞춰 세션을 재시작한다.
\end{enumerate}

이와 유사하게 프로세스가 타임아웃 또는 크래시로 종료되면 \code{\_recover\_after\_crash}가 프로세스 그룹을 회수하고 새 \code{REPLProcess}를 기동한 뒤 마지막 정상 노드의 \code{active\_lineage()}를 루트부터 결정론적으로 재실행하여 세션을 자동 복구한다. 추가로 \code{pickle}/\code{unpickle}(.olean 직렬화 연동), \code{format\_tactic\_input}(여러 줄 전술을 \code{(\textbackslash n  ...\textbackslash n)} 블록으로 묶어 REPL 단일 전술 호출에 적합하게 변환), \code{inspect\_symbol}(\code{\#check}/\code{\#print} 질의), \code{undo}, \code{reset}, \code{switch\_project} 메서드를 제공한다.

\subsection{\code{InfoTree} 재귀 추출 및 하이브리드 InfoView (\code{infoview.py}, \code{renderer.py})}
VS Code Lean 4 확장과 유사한 커서 연동 증명 검사 경험을 노트북 출력 셀에 구현하기 위해 \code{infoview.py}(1,260라인)는 REPL이 \code{infotree} 옵션으로 반환하는 재귀 JSON 트리를 순회한다. \code{extract\_infoview\_payload}는 트리의 \code{TacticInfo} 및 \code{TermInfo} 노드를 방문하여 각 스텝을 \code{TacticStepInfo}(\code{goalsBefore} $\rightarrow$ \code{goalsAfter} 및 \code{subgoals\_opened})와 \code{TermItemInfo}로 정규화한다. 이때 구조적 구두점이나 서식성 노드 등 실질적 상태 변화가 없는 노드는 \code{NON\_SUBSTANTIVE\_TACTIC\_KINDS} 집합으로 필터링하며, 트리 정보가 누락된 환경에서는 평면 \code{tactics} 리스트로 안전하게 폴백한다. 또한 비선형 산술 목표에서 \code{omega} 전술이 실패할 경우 \code{\_augment\_error\_message}가 한국어/영문 이중 안내 힌트를 부가하여 학습자 디버깅을 돕는다.

렌더링 아키텍처는 \textbf{SSR + CSR 하이브리드 위젯} 구조를 채택하였다. \code{\_render\_initial\_info\_pane\_html}이 초기 목표 상태를 서버 측에서 즉시 HTML로 생성하므로 자바스크립트가 비활성화된 정적 \code{.ipynb} 뷰어에서도 첫/마지막 상태가 표시된다. 이후 브라우저에서 \code{render\_infoview\_html}의 스크립트가 마운트되면 좌우 분할(\code{46\% / 1fr}, 뷰포트 720px 이하 반응형 전환, 각 패널 최대 높이 440px) UI가 활성화되며, \code{Prev}/\code{Next} 스텝 이동, \code{Before→After}/\code{After}/\code{Before} 보기 전환, \code{Expand} 토글, 그리고 \code{bindParentEditorCursor} 훅을 통해 부모 에디터(Web Studio의 \code{textarea}, Jupyter Classic의 CodeMirror 5 \code{.CodeMirror}, JupyterLab의 CodeMirror 6 \code{.cm-content})의 캐럿 오프셋에 대응하는 전술 스텝을 실시간 하이라이트한다. \code{renderer.py}의 \code{render\_outcome}은 \code{text/plain}과 \code{text/html} MIME을 동시 방출하되, 단순 선언 성공으로 출력이 비어 있고 InfoView 스텝도 없을 때는 \code{None}을 반환해 불필요한 공백 출력을 억제하며, turnstile 기호($\vdash$) 강조와 \code{all\_tactics}·원시 \code{InfoTree} 접힘 섹션을 제공한다.

\subsection{매직 명령, 구문 완결성 및 커널 통합 (\code{magics.py}, \code{kernel.py})}
\code{Lean4Kernel}(\code{ipykernel.kernelbase.Kernel} 상속)은 표준 Jupyter 메시지 핸들러를 구현한다.
\begin{itemize}
  \item \textbf{실행 (\code{do\_execute}):} 셸 요청의 \code{parent.metadata.cellId}를 읽어 세션 노드에 기록하며, 전술 실패 시에도 부분 \code{InfoViewPayload}를 추출해 오류 지점 직전의 목표 상태를 즉시 시각화한다.
  \item \textbf{매직 명령 (\code{magics.py}):} \code{MAGIC\_NAMES}에 정의된 17종 매직—상태 조회(\code{\%help}, \code{\%lsmagic}, \code{\%env}, \code{\%history}, \code{\%infoview}, \code{\%goals}), 전술 실행(\code{\%tactic}, \code{\%\%tactic}), 이력/스냅샷(\code{\%undo}, \code{\%reset}, \code{\%load}, \code{\%pickle}, \code{\%unpickle}), 프로젝트/원시 제어(\code{\%project}, \code{\%config}, \code{\%repl\_json}, \code{\%\%repl\_json})—을 지원한다.
  \item \textbf{완성 및 검사 (\code{do\_complete}, \code{do\_inspect}):} \code{unicode\_map.py}의 172개 역슬래시 유니코드 약어(\code{\textbackslash forall} $\rightarrow \forall$ 등) 확장을 최우선 처리한 뒤 매직 명령, 세션 내 선언된 사용자 식별자, 62개 \code{LEAN4\_KEYWORDS}, 56개 \code{LEAN4\_TACTICS} 순으로 후보를 제시한다. \code{do\_inspect}는 \code{inspect\_symbol}로 타입과 정의를 팝오버에 반환한다.
  \item \textbf{구문 완결성 (\code{do\_is\_complete}):} 중첩 주석 \code{/- ... -/}, 문자열 리터럴, 익명 생성자 쌍 $\langle\dots\rangle$ 균형 및 줄 끝 연접 접미사를 추적하여 다중 행 REPL 입력 대기 상태를 판별한다.
\end{itemize}

\section{Rocq (Coq) 커널 엔진 (\code{rocq\_kernel})}

\subsection{설계 결정 및 대안 비교}
\code{docs/designs/rocq\_kernel.md}~\cite{rocq_design_2026}에 문서화된 아키텍처 검토 과정에서 Rocq 백엔드를 구현하기 위한 세 가지 대안(\code{SerAPI}, \code{vscoq}/\code{coq-lsp}, 일괄 \code{coqc})이 평가되었다. \code{SerAPI}는 별도 OPAM 설치가 필요하고 최신 Rocq 버전 호환 유지가 중단되는 위험이 있었으며, \code{vscoq}/\code{coq-lsp}는 LSP 프레임 오버헤드와 배포 버전 파편화 문제가 존재하였다. 반면 배치 컴파일러 \code{coqc}는 셀 단위 증분 백트래킹과 문장별 목표 추적이 불가능하다. 이에 따라 본 프로젝트는 Coq/Rocq 기본 배포판에 내장된 \code{coqidetop.opt}/\code{coqidetop}의 네이티브 XML 프로토콜(\code{-main-channel stdfds})을 직접 구현하는 방안을 채택하였다.

\subsection{\code{coqidetop} XML 프로토콜 드라이버 (\code{xml\_protocol.py}, \code{driver.py})}
\code{RocqProcess}는 \code{find\_coqidetop\_binary} 및 \code{detect\_rocq\_version}으로 탐지한 \code{coqidetop} 프로세스를 \code{-main-channel stdfds} 플래그로 기동한다. XML 프레이밍은 \code{threading.Lock} 내부에서 \code{select.select}를 이용해 응답 종료 태그 \code{</value>}가 도착할 때까지 읽은 뒤, 비동기 피드백 XML 조각을 유실하지 않도록 추가로 \code{5\,ms} 후행 드레인(Trailing Drain)을 수행한다.

프로토콜 호환성과 파싱 안정성을 위해 \code{\_sanitize\_xml\_stream}은 XML 1.0 비호환 제어 바이트를 제거하고 \code{\&nbsp;} 엔티티를 \code{\&\#160;}으로 치환한다. \code{parse\_xml\_response}는 Coq/Rocq 세대 간 스키마 차이(\code{<state\_id>}와 \code{<stateid>}, \code{<hyp>}, \code{<richpp>})를 일관된 \code{RocqValue}, \code{RocqFeedback}, \code{RocqGoalsBundle}(\code{focused}, \code{background}, \code{shelved}, \code{admitted}) 데이터클래스로 파싱하며, 7종의 XML 원시 호출(\code{Init}, \code{Add}, \code{Edit\_at}, \code{Goal}, \code{Status}, \code{Query}, \code{Quit})과 \code{SIGINT} 인터럽트를 지원한다.

\subsection{어휘 기반 문장 토크나이저 및 셀 원자적 롤백 (\code{session.py})}
Lean 4가 블록 단위로 JSON 명령을 평가하는 것과 달리, Rocq는 마침표(\code{.})로 종결되는 문장(Vernacular/Tactic) 단위로 STM 상태 ID를 전진시킨다~\cite{barras2015asynchronous}. \code{tokenize\_sentences}는 순수 정규식 분할 시 발생하는 오분할을 막기 위해 다음 어휘 규칙을 상태 기계로 구현하여 \code{SentenceSpan} 목록을 생성한다.
\begin{itemize}
  \item 중첩 Coq 주석 \code{(* ... *)} 내부의 마침표 무시.
  \item 이중 따옴표 이스케이프 \code{""}를 포함한 문자열 리터럴 내부 마침표 무시.
  \item 범위 표기 줄임표 \code{..}와 식별자 접근자(모듈 경로 \code{Nat.add} 등) 구분: 마침표 직후 공백 또는 파일 끝(EOF)이 뒤따를 때만 문장 종결로 확정.
  \item 증명 구조 불릿과 중괄호(\code{-}, \code{+}, \code{*}, \code{\{}, \code{\}})를 독립 실행 단위로 분리하는 \code{\_split\_bullet\_subcommands}.
\end{itemize}

\begin{figure}[!t]
\centering
\begin{tikzpicture}[
  font=\scriptsize,
  nd/.style={draw, rounded corners=2pt, minimum width=1.15cm, minimum height=5.5mm, align=center},
  arr/.style={-{Stealth[length=1.8mm]}, semithick}
]
% Part (a) Lean 4 DAG Backtrack
\node[anchor=west, font=\scriptsize\bfseries] at (0, 2.35) {(a) \code{LeanSession}: \code{EnvNode} DAG 재선언 자동 백트래킹};
\node[nd, fill=blue!8] (e0) at (0.65, 1.65) {$e_0$\\{\tiny Root}};
\node[nd, fill=gray!15] (e1) at (2.35, 1.65) {$e_1$\\{\tiny \code{def f}}};
\node[nd, fill=gray!15] (e2) at (4.10, 1.65) {$e_2$\\{\tiny \code{lem}}};
\node[nd, fill=green!10] (e3) at (2.35, 0.70) {$e_3$\\{\tiny \code{def f'}}};
\node[nd, fill=green!10] (e4) at (4.10, 0.70) {$e_4$\\{\tiny replay \code{lem}}};

\draw[arr, dashed] (e0) -- (e1);
\draw[arr, dashed] (e1) -- (e2);
\draw[arr] (e0) |- node[pos=0.73, below, yshift=-1pt, font=\tiny]{재실행 분기} (e3);
\draw[arr] (e3) -- (e4);

% Part (b) Rocq STM Rollback
\node[anchor=west, font=\scriptsize\bfseries] at (0, -0.15) {(b) \code{RocqSession}: STM 셀 진입점 원자적 롤백 (\code{Edit\_at})};
\node[nd, fill=blue!8] (s0) at (0.65, -0.95) {$s_0$\\{\tiny Init}};
\node[nd, fill=blue!12] (s1) at (2.35, -0.95) {$s_1$\\{\tiny Cell Entry}};
\node[nd, fill=yellow!15] (s2) at (4.10, -0.95) {$s_2$\\{\tiny Sent. 1 OK}};
\node[nd, fill=red!14] (s3) at (5.85, -0.95) {$s_3$\\{\tiny Sent. 2 Err}};

\draw[arr] (s0) -- (s1);
\draw[arr] (s1) -- node[above, font=\tiny]{\code{Add}} (s2);
\draw[arr] (s2) -- node[above, font=\tiny]{\code{Add}} (s3);
\draw[arr, red!70!black, bend left=36] (s3.south) to node[below, font=\tiny]{\code{Edit\_at}($s_1$) 원자적 롤백 (\code{rollback\_whole\_cell\_on\_error})} (s1.south);
\end{tikzpicture}
\caption{두 증명기 커널의 상태 백트래킹 비교. (a) Lean 4는 불변 환경 스냅샷 DAG에서 충돌 조상의 \code{parent\_env}로 분기해 후손을 재실행하고, (b) Rocq는 셀 내부 부분 문장 실패 시 \code{Edit\_at}으로 셀 진입 STM 상태 ID($s_1$)까지 즉시 원상 복구한다.}
\label{fig:dag_rollback}
\end{figure}

\code{RocqSession}은 \code{STMNode} 계보와 \code{\_cell\_entry\_state} 매핑을 결합해 두 가지 안전장치를 제공한다(Fig.~\ref{fig:dag_rollback}(b)). 첫째, \code{auto\_backtrack}이 동일 \code{cell\_id} 또는 \code{extract\_declared\_name} 충돌을 감지하면 해당 진입 상태로 \code{Edit\_at}을 호출한 뒤 실행한다. 둘째, 하나의 셀 안에 여러 문장이 있을 때 후속 문장에서 오류가 발생하면 부분 적용 상태를 남기지 않고 \code{rollback\_whole\_cell\_on\_error=True} 정책에 따라 셀 실행 전 상태 $s_1$로 \code{Edit\_at} 원자적 롤백을 수행한다. 크래시 시에는 \code{\_recover\_after\_crash}가 \code{STMNode} 계보를 재실행하고, \code{inspect\_symbol}은 XML \code{Query} 명령을 통해 타입·정의 검색을 수행한다.

\subsection{Rocq InfoView 및 두 커널 간 아키텍처 비교}
\code{rocq\_kernel/infoview.py}(685라인)는 각 문장 실행 직후 \code{Goal} RPC로 수집한 \code{RocqGoalsBundle}을 \code{RocqStepInfo}로 변환한다. Lean InfoView와 동등한 split-pane 위젯(\code{Before→After}/\code{After}/\code{Before}, focused/background/shelved/admitted 섹션 구분, SSR 폴백)을 제공하며, 웹 스튜디오 프론트엔드와 직접 연동되는 \code{root.selectCursor} 메서드를 노출한다. \code{RocqKernel}은 8종의 매직(\code{\%help}, \code{\%lsmagic}, \code{\%goals}, \code{\%state}, \code{\%history}, \code{\%undo}, \code{\%reset}, \code{\%load})과 123개 유니코드, 66개 \code{ROCQ\_KEYWORDS}, 74개 \code{ROCQ\_TACTICS}를 지원한다. 두 증명기 커널의 기술적 차이는 Table~\ref{tab:lean_vs_rocq}에 요약하였다.

\begin{table}[!t]
\caption{Lean 4 커널과 Rocq 커널 아키텍처 상세 비교}
\label{tab:lean_vs_rocq}
\centering
\scriptsize
\begin{tabularx}{\columnwidth}{@{}l X X@{}}
\toprule
\textbf{비교 항목} & \textbf{\code{lean4\_kernel}} & \textbf{\code{rocq\_kernel}} \\
\midrule
하위 프로세스 & \code{repl} / \code{lake env repl} (\code{v4.35.0}) & \code{coqidetop.opt} / \code{coqidetop} \\
IPC 프로토콜 & JSON + \code{\textbackslash n\textbackslash n} stdin/stdout & XML (\code{</value>} + 5\,ms 드레인) \\
상태 추상화 & 불변 \code{EnvNode}(\code{env}), \code{ProofNode} & \code{STMNode}(\code{state\_id}), 셀 진입점 \\
구문 분할 & REPL 전체 파싱 + \code{InfoTree} & \code{tokenize\_sentences} + 불릿 분리 \\
재실행 복구 & \code{active\_lineage} 역추적·후손 재실행 & \code{Edit\_at} 후퇴 및 재실행 \\
부분 오류 처리 & 부분 \code{InfoView} 반환 + 힌트 & \code{Edit\_at} 셀 원자적 롤백 \\
크래시 복구 & \code{killpg} 후 계보 재실행 & \code{Quit}/\code{kill} 후 계보 재실행 \\
InfoView 계층 & \code{TacticInfo}/\code{TermInfo}, 1,260라인 & \code{RocqGoalsBundle} 4섹션, 685라인 \\
매직 / 어휘 맵 & 매직 17 · Uni 172 · KW 62 · Tac 56 & 매직 8 · Uni 123 · KW 66 · Tac 74 \\
\bottomrule
\end{tabularx}
\end{table}

\section{Web Proof Studio 및 Project-Agora 플랫폼 (\code{web.py}, \code{static/})}

\subsection{HTTPS 전용 보안 서버 및 인증 (\code{ProofNotebookHTTPServer})}
웹 서버 모듈 \code{web.py}(\code{src/\allowbreak lean4\_kernel/}, 3,395라인)는 \code{Threading\-HTTP\-Server} 기반의 웹 애플리케이션 서버(\code{Proof\-Notebook\-HTTP\-Server}, 기본 포트 \code{8890})를 제공한다. \code{resolve\_tls\_credentials}는 평문 HTTP 서비스를 허용하지 않고 TLS $\ge 1.2$ 기반 HTTPS만을 강제하며, 외부 인증서가 지정되지 않은 로컬 개발 환경에서는 \code{openssl}을 호출해 EC P-256 자체 서명 인증서를 자동 생성한다. \code{ProofStudioApp}은 연구자 JWT 토큰을 SHA-256 해시 키로 120초 TTL 캐싱하여 반복 인증 부하를 줄이면서 사용자별 워크스페이스 격리를 유지한다.

\subsection{멀티 커널 LRU 풀 (\code{ProofKernelPool})}
표준 Jupyter 서버가 노트북 파일 전체를 하나의 고정 커널에 묶는 것과 달리, \code{ProofKernelPool}은 단일 워크스페이스 내부에서 셀 언어(\code{\_normalize\_language}: \code{lean4}, \code{rocq}/\code{coq}, \code{python})에 따라 독립된 ZMQ \code{KernelManager}를 라우팅한다.
\begin{itemize}
  \item \textbf{키 및 용량 제한:} 커널 풀 키는 \code{user\_id::notebook\_id::lang} 형식을 따르며, 최대 동시 실행 커널 수를 $k_{\max} = 16$으로 제한하는 LRU(Least Recently Used) 캐시로 메모리 고갈을 방지한다.
  \item \textbf{인메모리 커널 스펙 합성:} 호스트에 \code{jupyter kernelspec install}이 사전에 수행되지 않은 환경에서도 모듈 경로를 이용해 임시 커널 스펙을 즉석 합성하여 커널을 기동할 수 있다.
  \item \textbf{출력 트렁케이션 가드:} 대규모 로그나 시각화로 브라우저·Redis가 포화되지 않도록 셀당 출력 항목 수 \code{MAX\_OUTPUT\_ITEMS=64}, 스트림 문자 수 \code{MAX\_STREAM\_CHARS=128\,KiB}, HTML 출력 \code{MAX\_HTML\_CHARS=512\,KiB}, 이미지 \code{MAX\_IMAGE\_CHARS=2\,MiB} 상한을 강제한다.
\end{itemize}
또한 \code{ProofKernelPool.complete}는 활성 ZMQ 커널 셸 채널 우선 조회 및 빠른 유니코드/키워드 로컬 완성을 제공하며, \code{inspect\_symbol}은 커널 \code{inspect\_request}를 통해 심볼 시그니처를 조회한다.

\subsection{제로 의존성 Redis 및 HOT/COLD 생명주기 (\code{WorkspaceLifecycleManager})}
\code{AgoraRedisClient}는 \code{redis-py} 같은 외부 패키지 의존성 없이 Python 표준 \code{socket}과 \code{ssl}만으로 Redis RESP2 프로토콜(TCP 및 TLS, ACL \code{AUTH} 지원, 기본 엔드포인트 \code{172.20.0.2:6379}, 기본 TTL \code{3600\,s})을 직접 구현하였다. 또한 테스트나 오프라인 단독 실행 환경에서 Redis 접속이 불가할 경우 동일한 API 시맨틱을 유지하는 스레드 안전 인메모리 폴백 저장소로 자동 전환된다.

\begin{figure}[!t]
\centering
\begin{tikzpicture}[
  font=\scriptsize,
  st/.style={draw, rounded corners=3pt, text width=4.4cm, align=center, inner sep=4pt, minimum height=0.88cm},
  arr/.style={-{Stealth[length=2mm]}, semithick}
]
\node[st, fill=gray!10] (cold) {\textbf{COLD 상태} (디스크 영속화)\\\code{users/<uid>/notebooks} JSON 저장 (활성 커널 없음)};
\node[st, fill=orange!12, below=0.80cm of cold] (hot) {\textbf{HOT 상태} (Redis + 실시간 협업)\\\code{proof:workspace:<id>:\{doc,meta,clients\}}\\\code{cacheGeneration++} · LRU 커널 · SSE 활성화};

\draw[arr] ([xshift=-1.1cm]cold.south) to[bend right=20] node[left, align=right, font=\scriptsize]{\textbf{Mount / Execute}\\디스크 로드 + 3600\,s TTL} ([xshift=-1.1cm]hot.north);
\draw[arr] ([xshift=1.1cm]hot.north) to[bend right=20] node[right, align=left, font=\scriptsize]{\textbf{Unmount / Evict}\\디스크 플러시 + 커널 종료} ([xshift=1.1cm]cold.south);
\draw[arr] ([xshift=0.9cm]hot.south) .. controls +(0.15,-0.40) and +(-0.15,-0.40) .. node[below, font=\scriptsize]{SSE 브로드캐스트 · Presence (45\,s TTL)} ([xshift=-0.9cm]hot.south);
\end{tikzpicture}
\caption{\code{WorkspaceLifecycleManager}의 Redis \code{COLD $\leftrightarrow$ HOT} 상태 전이 모델. 워크스페이스 마운트 시 \code{cacheGeneration}이 증가하며 언마운트 시 안전하게 디스크로 플러시되고 해당 ZMQ 커널이 회수된다.}
\label{fig:lifecycle}
\end{figure}

Fig.~\ref{fig:lifecycle}과 같이 \code{WorkspaceLifecycleManager}는 \code{proof:workspace:<id>:\{doc,meta,clients\}} 네임스페이스 키를 통해 워크스페이스를 \code{COLD}(디스크 저장)에서 \code{HOT}(Redis 캐시 마운트) 상태로 승격한다. 동시 편집 충돌이나 캐시 무효화를 식별하기 위해 단조 증가하는 \code{cacheGeneration} 카운터를 유지하며, 사용자가 워크스페이스 탭을 언마운트하면 최신 문서를 디스크 \code{NotebookStore}(\code{users/\allowbreak<uid>/\allowbreak notebooks} 및 \code{shared\_notebooks})에 플러시하고 연관 ZMQ 커널들을 즉시 종료한다. \code{NotebookStore}는 표준 \code{nbformat 4.5} 기반 \code{to\_ipynb}/\code{from\_ipynb} 직렬화와 \code{to\_lean\_source}(Markdown 및 Python 셀을 Lean 주석으로 감싸 단일 \code{.lean} 모듈로 내보내기)를 지원한다.

\subsection{실시간 SSE 협업, Project-Agora 브리지 및 프론트엔드}
다중 연구자의 동시 편집을 위해 \code{RealtimeNotebookHub}는 Server-Sent Events(SSE) 기반 pub/sub 채널과 45초 TTL 기반 현재 접속자(Presence) 추적을 제공한다. \code{AgoraBridge}는 \code{\_read\_agora\_be\_env}와 \code{\_detect\_agora\_ca\_cert}를 통해 동일 호스트 상의 IRADE/Agora 백엔드 \code{.env} 설정과 로컬 루트 CA 인증서를 자동 탐색하고, 이메일·MFA 로그인, JWT 연동 및 협업 캔버스(\code{notebook}/\code{text} 카드) 양방향 가져오기·내보내기(\code{httpx} 또는 선택적 \code{agora\_canvas.AgoraClient})를 수행한다.

프론트엔드(\code{index.html}, \code{app.js}, \code{styles.css})는 VS Code의 Activity Bar, Explorer(Redis HOT/All Workspaces), Project-Agora Canvas Bridge, Symbols \& Magics 팔레트, 다중 탭 에디터 및 상태 바 레이아웃을 재현한다. 특히 \code{executeCellById}가 \code{POST /api/workspaces/<id>/execute}로부터 받은 InfoView HTML은 \code{mountInteractiveHtml}을 통해 동적 \code{<script>} 태그를 재주입·실행함으로써 커서 동기화 위젯이 브라우저 DOM 상에서 즉시 동작하도록 보장한다. 또한 \code{EventSource}(\code{ensureWorkspaceEventStream}), \code{\textbackslash 기호 + Tab} 유니코드 변환(\code{tryExpandBackslashUnicode}), \code{Shift+Tab} 심볼 팝오버(\code{triggerInspect}) 및 클라이언트 \LaTeX/Markdown 렌더러(\code{renderMarkdownHtml})를 제공한다.

\section{검증 체계, 활용 예제 및 통합 현황 평가}

\subsection{CLI 진단 도구 및 형식 증명 예제}
설치 및 환경 자가 진단은 양 커널의 \code{install.py}가 담당한다. \code{lean4-kernel} CLI는 \code{install}, \code{build-repl}, \code{doctor}(실제 \code{\#eval 40 + 2} 실행 후 출력 \code{"42"} 검증), \code{uninstall}, \code{serve} 서브커맨드를 제공하며, \code{rocq-kernel} CLI는 \code{install}과 \code{doctor}(\code{Check nat.} 피드백 내 \code{"Set"} 포함 여부 검증)를 제공한다.

저장소의 \code{examples/} 및 \code{notebooks/}에는 두 증명기의 상호작용을 시연하는 검증된 모듈이 포함되어 있다.
\begin{itemize}
  \item \code{examples/Arithmetic.lean}(23라인): 재귀 함수 \code{gaussSum}과 가우스 합 항등식 \code{gaussSum\_eq} ($2 \sum_{k=0}^{n} k = n(n+1)$)를 \code{sorry} 없이 \code{calc} 블록으로 증명하고 \code{\#print axioms gaussSum\_eq}로 공리 건전성을 검증한다.
  \item \code{examples/Arithmetic.v}(43라인): \code{double}, \code{gauss\_sum} 정의와 함께 \code{plus\_O\_n}, 불릿(\code{-}) 기반 귀납법 \code{add\_0\_r} 및 \code{add\_assoc}를 \code{Admitted} 없이 증명하고 \code{Check add\_assoc.}을 실행한다.
  \item \code{notebooks/lean4\_demo.ipynb}(15셀) 및 \code{notebooks/rocq\_demo.ipynb}(9셀): \code{\%load}, \code{\%tactic}, \code{\%undo}, InfoView 및 백트래킹 시나리오를 수록한 데모 노트북.
\end{itemize}

\subsection{자동화 테스트 스위트 및 샌드박스 검증 결과}
테스트 스위트는 총 6개 파일, 30개 테스트 케이스(1,578라인)로 구성된다. 작업 공간이 기존 경로(\code{proof-notebook})에서 \code{notebook-kernel}로 이동한 환경에서는 \code{.venv/bin/pytest} 래퍼의 shebang 경로가 이전 디렉터리를 가리키므로 \code{PYTHONPATH=src .venv/bin/python3 -m pytest -v} 형태로 실행하였다. 현 세션에서 직접 실행·관측한 전체 결과는 \textbf{21 passed, 2 skipped, 7 failed (7.22\,s)}이며, 모듈별 세부 내역은 Table~\ref{tab:test_results}와 같다.

\begin{table}[!t]
\caption{자동화 테스트 스위트 실행 결과 (\code{PYTHONPATH=src .venv/bin/python3 -m \mbox{pytest -v}})}
\label{tab:test_results}
\centering
\scriptsize
\begin{tabularx}{\columnwidth}{@{}l@{\hspace{5pt}} c@{\hspace{4pt}} c@{\hspace{4pt}} c@{\hspace{6pt}} X@{}}
\toprule
\textbf{테스트 파일} & \textbf{통과} & \textbf{스킵} & \textbf{실패} & \textbf{범위 및 실패/스킵 사유} \\
\midrule
\code{test\_repl\_protocol.py} & 10 & 0 & 0 & Lean JSON 파싱, C stdout 분리, DAG, InfoView SSR (오프라인) \\
\code{test\_rocq\_protocol.py} & 6 & 0 & 0 & Rocq XML 파싱, \code{tokenize\_sentences}, \code{MockRocqDriver}, InfoView \\
\code{test\_rocq\_session.py} & 2 & 0 & 0 & 완성 검사 및 실제 \code{coqidetop} 라이브 세션 통합 검증 \\
\code{test\_session\_and\_magics.py} & 2 & 1 & 0 & 유니코드·완결성 통과; 라이브 Lean 세션은 \code{PATH} 내 \code{lean} 부재로 스킵 \\
\code{test\_notebook\_e2e.py} & 0 & 1 & 0 & \code{nbclient} E2E 실행 (\code{lean} 부재로 스킵) \\
\code{test\_web\_platform.py} & 1 & 0 & 7 & \code{NotebookStore} 단위 통과; HTTPS 테스트 7건은 샌드박스 프록시가 루프백(\code{127.0.0.1})을 차단하여 \code{httpx.ProxyError: 400 Bad Request} 발생 \\
\midrule
\textbf{합계 (총 30건)} & \textbf{21} & \textbf{2} & \textbf{7} & 오프라인(16/16) 및 라이브 Rocq(2/2) 전량 통과 \\
\bottomrule
\end{tabularx}
\end{table}

결과가 보여주듯 오프라인 프로토콜·토크나이저·세션 백트래킹 단위 테스트(16건)와 시스템에 설치된 \code{coqidetop} 기반의 라이브 Rocq 세션 통합 테스트(2건)는 오류 없이 통과하였다. 반면 웹 플랫폼 HTTP 통합 테스트 7건은 애플리케이션 로직 결함이 아닌 샌드박스 아웃바운드/루프백 프록시 환경에서 \code{httpx} 클라이언트가 \code{127.0.0.1} 트래픽을 프록시로 우회시키며 발생한 \code{ProxyError}에 기인한다.

\subsection{관찰된 통합 한계 및 향후 개선 과제}
코드 감사와 실행 검증을 통해 확인된 현재 아키텍처의 구체적 개선 과제는 다음과 같다.
\begin{enumerate}
  \item \textbf{Web Proof Studio UI 및 저장소의 Rocq 노출 보완:} 백엔드 \code{ProofKernelPool}은 이미 \code{"rocq"}/\code{"coq"} 언어를 \code{rocq\_kernel}로 정상 라우팅하나, 프론트엔드(\code{index.html}, \code{app.js})의 셀 추가 버튼·상태 배지(\code{Lean 4} / \code{Python} 이분 표시), \code{ProofKernelPool.complete}의 정적 키워드 폴백(\code{LEAN\_KEYWORDS} 고정 참조), 그리고 \code{NotebookStore}의 \code{list\_notebooks} 통계 및 \code{to\_lean\_source} 직렬화는 아직 Lean 4와 Python 중심으로 구현되어 있어 UI 차원의 Rocq 셀 선택자 및 \code{.v} 내보내기 확장이 필요하다.
  \item \textbf{커널 간 공통 추상화 추출:} \code{lean4\_kernel/infoview.py}(1,260라인)과 \code{rocq\_kernel/infoview.py}(685라인), 그리고 두 \code{session.py}의 DAG 계보 추적·크래시 복구 코드는 상당 부분 대칭적 구조를 공유하므로, 증명기 중립적 공통 패키지(\code{proof\_kernel\_core})로 추출하면 유지보수 중복을 크게 줄일 수 있다.
  \item \textbf{프록시·가상환경 경로 독립성 강화:} \code{test\_web\_platform.py} 등의 테스트 \code{httpx} 클라이언트에 \code{trust\_env=False}를 지정해 프록시 환경의 루프백 간섭을 배제하고, \code{.venv} 재배치 시에도 스크립트 shebang이 깨지지 않도록 editable 재설치를 표준화할 필요가 있다.
\end{enumerate}

\section{결론}
본 기술 보고서에서는 Lean 4, Rocq(Coq), 그리고 Python 3를 하나의 대화형 연구 워크벤치로 통합한 \textit{Proof-Notebook}의 전체 아키텍처와 구현 세부를 IEEEtran 형식으로 기록하였다. Lean 4의 불변 \code{EnvNode} DAG 및 재귀 \code{InfoTree} 추출, Rocq의 \code{coqidetop} XML 프로토콜 드라이버 및 셀 원자적 \code{Edit\_at} 롤백, 그리고 TLS 1.2+ 웹 스튜디오·제로 의존성 RESP2 Redis HOT/COLD 생명주기·SSE 실시간 동기화는 정리 증명기의 엄밀성과 노트북 환경의 대화형 편의성을 양립시키는 견고한 기반을 형성한다. 오프라인 프로토콜 검증과 라이브 Rocq 세션 통합을 통해 핵심 상태 기계의 신뢰성이 확인되었으며, 향후 웹 UI 내 Rocq 셀 일급 지원 및 공통 InfoView/세션 추상화가 완료되면 진정한 의미의 범용 다중 증명기 협업 R\&D 플랫폼으로 완성될 것이다.

\bibliographystyle{IEEEtran}
\bibliography{references}

\end{document}
